\scanpage{9}
\diagram{\scriptsize\begin{tikzcd}[ampersand replacement=\&,column sep=.8em]
0 \arrow[r] \& \deleted{\unsure{\mathbb Z(G)}}\,0 \arrow[r] \& \mathbb Z^I \arrow[r,dashed] \&
\underset{\text{Teichmüller spécial}}{TS(C,D,\varphi)} \arrow[r,two heads,"\psi"] \&
\underset{\substack{\text{Teichmüller}\\\text{à trous}}}{T(C,D,\varphi)} \arrow[r] \&
\underset{\substack{\text{groupe}\\\text{symétrique}}}{\unsure{\mathfrak S_I|\mathfrak G_I}} \arrow[r] \& 1
\end{tikzcd}}
% slanted
\margin{nul sauf pour la sphère à 1 ou 2 trous}
Les $(v,\alpha)\in TS$ dans le noyau de $\psi:TS\to T$, i.e. tels que $\widetilde v=\mathrm{id}$, sont représentables par des $(v,\alpha)$ avec $v=\mathrm{id}$, donc
\[
h_i=\mathrm{int}\,(\alpha(i))\,h_i
\]
Sauf dans le cas de la sphère à $1$ trou, \deleted{\unsure{on|on a}} le centralisateur de $h_i$ est le groupe qu'il engendre \uline{et} ce groupe est libre de rang $1$, donc on aura alors $\deleted{\unsure{\in}}$
\[
\alpha(i)=h_i^{\lambda_i}\qquad\lambda_i\in\mathbb Z\ \text{bien déterminé}
\]
On aura
\[
(\mathrm{id}_G,(h_i^{\lambda_i}))\sim(\mathrm{id}_G,(h_i^{\lambda'_i}))
\]
sss $\exists\mu\in Z(G)$ tel que
\[
 h_i^{\lambda'_i}=\mu\,h_i^{\lambda_i}
\]
Dans tous les cas sauf la sphère à $2$ trous (cylindre) $Z(G)=0$, donc la condition précédente est $(\lambda_i)=(\lambda'_i)$, \inserted{i.e. $\mathrm{Ker}\,\psi\simeq\mathbb Z^I$} Dans le cas de la sphère à $2$ trous, $Z(G)=G\underset{?}{\simeq}\mathbb Z$ \deleted{la condition est que} \inserted{$(h_1,h_2)=(+1,-1)$}
\[
\deleted{\mu=\unsure{h_i|\ell_i}\quad h_i^{\lambda'_i}\,/\,h_i^{\lambda_i}\quad\unsure{\text{i.e. de}}}
\qquad\deleted{\unsure{\text{or}}\quad h_1^{\lambda_1},h_2^{\lambda_2}}
\]
les conditions sont alors (notation additive de $\mathbb Z$)
\[
\begin{aligned}
\lambda'_1&=\mu+\lambda_1\\
-\lambda'_2&=\mu-\lambda_2\qquad\text{i.e.}\quad\lambda'_2=-\mu+\lambda_2
\end{aligned}
\]
i.e. $\mathrm{Ker}\,\deleted{\unsure{\psi}}\,(\mathbb Z^I\to\mathrm{Ker}\,\psi)\simeq\mathrm{Ker}\,(\mathbb Z^I\xrightarrow[\text{opération somme}]{}\mathbb Z)$

donc \qquad $\mathrm{Ker}\,\psi\simeq\mathbb Z$ \quad (canoniquement)

Enfin dans le cas de la sphère à $1$ trou, $G=0$ donc $\alpha=0$, donc $TS(C,D,\varphi)=0$

On trouve d'autre part, trivialement,
\[
\pi_1(\underline{\mathrm{End}}_{D\text{ fixé}}(C,D,\varphi))=1
\]
% archive sheet number: 8
